% ============================================================
\chapter{Introduction}
\label{chap:intro}
% ============================================================

\section{Context and motivation}
\label{sec:motivation}

Portfolio optimisation divides capital across assets so as to maximise return
for a given level of risk. Since Markowitz's mean--variance model
\cite{markowitz1952portfolio}, modern portfolio theory has decomposed an asset's
risk into a \emph{systematic} part, shared with the broader market and priced
through factor models such as the Capital Asset Pricing Model
\cite{sharpe1964capm} and Arbitrage Pricing Theory \cite{ross1976apt}, and an
\emph{idiosyncratic} part that diversification across assets reduces. These two
components ordinarily trade off: reducing factor exposure tends to concentrate a
portfolio, while diversifying across assets reduces variance but tends to
increase factor exposure.

Almost all of this machinery rests on the covariance, or equivalently the
correlation, matrix of returns. Correlation has three well-known weaknesses.
It is \emph{undirected} (it imposes no temporal ordering between assets'
movements); it is \emph{static} (it does not track how market structure shifts
during crises); and it can be \emph{spurious} (two variables may co-move with no
causal relationship \cite{simon1954spurious}). The last point is the most
damaging for factor investing: a variable that merely co-moves with equities,
such as an equity index, is statistically indistinguishable from one that
causally drives them, which manifests as overfitting and false factor
discoveries \cite{lopezdeprado2023causal}.

Causal discovery learns directed relationships from observational data and so, in
principle, separates genuine drivers from incidental co-movement. The field has
recently become tractable at scale: the score-based algorithm NOTEARS
\cite{zheng2018notears} and its temporal extension DYNOTEARS
\cite{pamfil2020dynotears} recast structure learning as continuous optimisation,
and a small but growing literature now applies these methods to equity markets
\cite{howard2025causal,tang2024trading}.

This project applies causal discovery to a recent allocation framework,
Hierarchical Sensitivity Parity (HSP) \cite{rodriguez2023hsp}. HSP first selects
the common \emph{drivers} of a set of assets (government bonds, commodity
futures, macroeconomic indicators, FX pairs, and so on) and then clusters assets
by the similarity of their \emph{sensitivities} to those drivers. Its appeal is
theoretical --- via the commonality principle and the Reichenbach common-cause
principle \cite{rodriguez2025causal,hitchcock2020reichenbach} it claims to escape
the systematic/idiosyncratic diversification trade-off --- yet it selects its
drivers by \emph{cumulative correlation}, the very quantity its causal motivation
was meant to move beyond. The central question of this dissertation is whether
replacing that correlation-based selection with genuine causal structure
produces more robust portfolios, particularly around the regime changes where
correlation-based methods are most vulnerable. I additionally test a
\emph{closed feedback loop} in which realised out-of-sample performance updates a
driver-utility score that shapes the next period's selection. I call HSP with
causally-selected drivers \emph{Causal-HSP}, and the version with feedback
\emph{closed-loop Causal-HSP}.

\section{The scientific question}
\label{sec:question}

The motivation above is sharpened into a single, precisely-stated question that
the rest of this dissertation answers in depth and returns to throughout.

\begin{quote}
\noindent\fbox{\parbox{\dimexpr\linewidth-2\fboxsep-2\fboxrule\relax}{%
\textbf{Q.}\ \emph{For a long-only, monthly-rebalanced portfolio of the
${\sim}100$ largest S\&P-500 constituents over 2007--2024, does choosing
Hierarchical Sensitivity Parity's common drivers by temporal causal discovery
(DYNOTEARS), rather than by cumulative correlation, deliver a higher
out-of-sample, net-of-cost, risk-adjusted return --- and does any such advantage
survive correction for the multiplicity of strategies tried and the
non-normality of returns?}}}
\end{quote}

Every term is operationalised and held fixed for the whole study:
\begin{itemize}[itemsep=1pt, topsep=3pt]
  \item \emph{Out-of-sample}: a walk-forward backtest --- drivers, sensitivities
    and weights at each rebalance use only data up to that date.
  \item \emph{Net of cost}: returns are charged $5$\,bps per unit of one-way
    turnover (swept over $\{0,5,10,20\}$\,bps).
  \item \emph{Risk-adjusted return}: primarily the annualised Sharpe ratio, but
    --- because daily returns here have excess kurtosis ${\approx}\,\rsKurtosis\
    ---$ also the Sortino and Calmar ratios and, decisively, the
    distribution-aware \emph{Probabilistic} and \emph{Deflated} Sharpe ratios
    \cite{bailey2012psr,bailey2014deflated} (Section~\ref{sec:method},
    Section~\ref{sec:believe-rigour}).
  \item \emph{Multiplicity}: the \rsNtrials\ distinct configurations actually
    evaluated (variants $\times$ discovery methods $\times$ windows $\times$ $K$
    $\times$ feedback grid), against which the Deflated Sharpe, Reality Check,
    SPA and Model Confidence Set deflate.
\end{itemize}

\noindent Q decomposes into five sub-questions, each answered in its own section
of Chapter~\ref{chap:evaluation} and each tied to an achievement:
\begin{enumerate}[label=\textbf{Q\arabic*}, itemsep=1pt, topsep=3pt]
  \item Does causal driver selection beat cumulative-correlation selection on net
    out-of-sample risk-adjusted return? \emph{(the treatment; A1/A2,
    Section~\ref{sec:causal-vs-corr})}
  \item Are exogenous drivers necessary at all, or does asset-only causal
    structure suffice? \emph{(necessity; A2, Section~\ref{sec:v0prime-eval})}
  \item Is any edge robust to the lookback window and to the number of drivers
    $K$? \emph{(robustness; A1/A3, Section~\ref{sec:ksens})}
  \item Can a monthly performance-feedback loop be learned from at all?
    \emph{(feedback; A3, Section~\ref{sec:loop-eval})}
  \item Does any edge survive correction for multiplicity and non-normality?
    \emph{(rigour; A3, Section~\ref{sec:believe-rigour})}
\end{enumerate}

\section{Challenges, novel contributions, and quantified benefits}
\label{sec:challenges}

The project is framed around three technical challenges, the contributions that
address them, and the quantified benefits actually measured. Each challenge maps
onto Q: \textbf{C1} is the \emph{treatment} in Q (causal vs correlation
selection); \textbf{C2} is the \emph{mechanism}, and through the asset-only
control the \emph{necessity}, of using directional structure; \textbf{C3} is the
\emph{rigour and feedback} on which Q5 and Q4 turn.

\begin{enumerate}
  \item \textbf{Key technical challenges.}
  \begin{enumerate}[label=\textbf{C\arabic*}]
    \item \emph{Causal driver identification at portfolio scale.} HSP selects
      drivers by cumulative correlation, conflating co-movement with causation.
      Identifying genuinely causal drivers requires causal discovery over a large
      joint panel of drivers and assets, a directional prior (drivers cause
      assets, not the reverse), and enough computational tractability to repeat
      the discovery at every rebalance --- non-trivial, as many causal-discovery
      algorithms are expensive and not readily parallelisable.
    \item \emph{Using directional causal structure inside a symmetric allocator.}
      Hierarchical risk-parity methods cluster on a symmetric distance matrix, but
      causal graphs are directional. Naively symmetrising the causal adjacency
      would discard the very directional information that motivates the approach.
    \item \emph{Closed-loop feedback.} I hypothesise that letting realised
      performance feed back into selection improves robustness across regime
      changes. Such feedback is not part of HSP, so the framework must be extended
      and the hypothesis tested rather than assumed.
  \end{enumerate}

  \item \textbf{Contributions addressing each challenge.}
  \begin{enumerate}[label=\textbf{A\arabic*}]
    \item A \emph{causal factor discovery pipeline} (Chapter~\ref{chap:discovery})
      that runs temporal causal discovery (DYNOTEARS, with VARLiNGAM as a
      comparator, and a reduced-scope NTS-NOTEARS probe) over the joint panel,
      enforces the directional prior by forbidding asset-to-driver edges,
      calibrates the number of drivers from the data, and is made reproducible and
      tractable by a content-keyed discovery cache.
    \item \emph{Causal-HSP in sensitivity space} (Chapter~\ref{chap:sensitivity}):
      each asset's sensitivities to the selected drivers form a Euclidean
      embedding whose distance matrix is symmetric \emph{by construction}, so
      directional causal information drives the clustering with no symmetrisation
      of the graph. This also enables an \emph{asset-only} control (V0$'$) that
      converts an asset--asset causal graph directly into weights.
    \item A \emph{closed feedback loop} (Chapter~\ref{chap:evaluation}) that
      attributes realised performance back to the drivers that shaped the
      portfolio and blends this utility with causal evidence in the next
      selection, evaluated across two lookback windows and a hyperparameter grid.
  \end{enumerate}

  \item \textbf{Quantified benefits actually measured} (full detail in
    Chapter~\ref{chap:evaluation}; net of 5\,bps costs, 215 monthly rebalances,
    2007--2024, significance by stationary block bootstrap
    \cite{politis1994stationary}). The results are reported here as measured,
    including where the method does \emph{not} help.
  \begin{itemize}[itemsep=1pt, topsep=3pt]
    \item At the \textbf{one-year} window causal selection beats the
      cumulative-correlation baseline (V0, Sharpe $0.371$): DYNOTEARS reaches
      $0.381$ (a $+0.010$ edge that is \emph{not} statistically significant,
      $p=0.27$) and VARLiNGAM reaches $0.399$ ($+0.028$, $p=0.004$).
    \item The single most effective variant is \textbf{asset-only Causal-HRP}
      (V0$'$, no exogenous drivers): Sharpe $0.403$, significantly above both the
      correlation baseline ($p<0.001$) \emph{and} the driver-based method
      ($p=0.01$). On this universe, asset--asset causal structure alone
      outperforms adding the exogenous-driver machinery.
    \item These edges are \textbf{not robust}. At the \textbf{two-year} window they
      vanish (every variant converges to Sharpe $\approx 0.372$, and VARLiNGAM
      falls below the baseline), and the V1$-$V0 edge changes sign with the number
      of drivers $K$.
    \item The \textbf{closed loop is inert}: across the entire $\alpha,\gamma$
      grid the feedback never changes the selected driver set, so V2 reproduces
      open-loop V1 exactly. The feedback hypothesis (C3) is answered in the
      negative.
    \item Qualitatively, the causal selector rotates across $\approx 26$ distinct
      drivers that track the macro regime, whereas the correlation selector locks
      onto the same $\approx 10$--$14$ drivers (international-equity ETFs and the
      rates curve) at every rebalance --- direct evidence of the cum-corr
      tautology the project set out to avoid.
    \item The closest prior work, Howard et al.\ \cite{howard2025causal}, applies
      DYNOTEARS to the S\&P~500 but never sets portfolio weights from the causal
      graph; converting causal structure directly into allocation (and the V0$'$
      result above) is the novel element here.
  \end{itemize}
\end{enumerate}

The honest one-line summary is therefore: \emph{causal structure delivers a
small, window-dependent improvement over correlation for this framework ---
clearest, and statistically significant, only as asset-only structure at the
short window --- and the closed-loop extension delivers nothing.} Establishing
this carefully, and reproducibly, is the contribution.

\section{Report structure}
\label{sec:structure}

Chapter~\ref{chap:background} reviews portfolio optimisation, causal discovery,
and their intersection, and positions the work against the literature.
Chapters~\ref{chap:discovery}--\ref{chap:evaluation} present the three
achievements: A1, the causal factor discovery pipeline (addressing C1);
A2, Causal-HSP in sensitivity space (C2); and A3 with the full empirical
evaluation, the closed loop and all comparative, regime-conditional and
robustness results (C3). Chapter~\ref{chap:conclusion} draws objective
conclusions, states the limitations of a causal (and causal-plus-asset) approach
plainly, and sets out future work.
